% TOP_PROBLEM: {"id": 20001163, "title": "Finite-field factorization and sparse-polynomial root testing", "category_id": 15, "category": {"id": 15, "name": "computer_science", "display_name": "Computer Science", "description": "Computational complexity, algorithms, and theoretical CS.", "slug": "computer-science", "order_index": 15, "created_at": "2026-07-31T15:26:25.671Z"}, "status": "open"}
% FIRST_POSTED: 2026-10-06
% ATTEMPT_STATUS: unresolved
% OPENAI_CLAIM: {"schema_version":1,"source_repo":"https://github.com/openai/math","source_commit":"adc7f1241b42e322a6451854ab7e4b4c146bf78a","release_date":"2026-10-06","match":"partial","resolution":"partial","independently_reviewed":false,"second_pass":"n/a","families":[{"family":"142","title":"Deterministic polynomial factorization over prime fields","subject":"Theoretical computer science","manuscripts":[{"title":"Deterministic Polynomial Factorization over Prime Fields","pdf_path":"preprints/Deterministic-Polynomial-Factorization-over-Prime-Fields-October-4-2026/Deterministic-Polynomial-Factorization-over-Prime-Fields.pdf","readme_path":"preprints/Deterministic-Polynomial-Factorization-over-Prime-Fields-October-4-2026/README.md","date":"2026-10-04","theorem_ref":"thm:factorization (sections/00-introduction.tex:27)","statement_sources":[{"path":"preprints/Deterministic-Polynomial-Factorization-over-Prime-Fields-October-4-2026/build/sections/00-introduction.tex","line":27,"label":"thm:factorization"}]}],"lean":null,"reasoning_trace":null}]}
\documentclass[11pt]{article}
\usepackage[margin=25mm]{geometry}
\usepackage{fontspec}
\setmainfont{Latin Modern Roman}
\usepackage{amsmath,amssymb,mathtools}
\usepackage{unicode-math}
\setmathfont{Latin Modern Math}
\usepackage{enumitem}
\usepackage{xurl,hyperref,bookmark}
\hypersetup{colorlinks=true,urlcolor=blue}
\setcounter{secnumdepth}{0}
\setlength{\parindent}{0pt}
\setlength{\parskip}{5pt}
\setlength{\emergencystretch}{3em}
\begin{document}
\section*{Finite-field factorization and sparse-polynomial root testing}

\subsection{OpenAI's claim}
OpenAI presents the following claims relevant to this question, with manuscript numbers following paper order in Sources. In manuscript 1 (family 142; 2026-10-04), the asserted result is that dense polynomials over binary-specified prime fields have uniform deterministic complete factorization in bit time polynomial in degree and $\log p$, using the cited Hecke zero-free input.

\subsection{Mathematical statement of the claimed result}
\paragraph{Manuscript 1.}
For any binary-encoded prime $p$ and nonzero densely encoded $f\in\mathbb F_p[x]$ of degree $n$, a uniform deterministic algorithm is claimed to output
\[
 f=c\prod_i g_i^{e_i}
\]
with $c$ the leading coefficient, distinct monic irreducible $g_i$ and positive integer multiplicities, in $O(((n+1)\lceil\log_2p\rceil)^{10^{12}})$ bit operations. Constants have an empty factor list. The computation is unconditional and deterministic; it receives neither an integer-factorization service nor a primitive-root oracle.

The separate algebraic reduction sets $L=\lceil\log_2p\rceil$, $B=20+(n+1)(L+1)$. If $p>B^{200000}$ and $n\ge2$, supply for every prime $q\le n$ a prime $\ell_q\notin\{2,3,p,q\}$ with
\[
 \ell_q\equiv1\pmod{12q},\qquad p^{(\ell_q-1)/q}\not\equiv1\pmod{\ell_q}.
\]
That reduction has bit cost $O((B+E)^C)$ for absolute $C$, where $E=\max(2,\max_q\ell_q)$, and $E=2$ when no such primes are required. This bound depends on the numerical size of the auxiliary primes; the full theorem additionally constructs sufficiently small primes.

\subsection{Outline of the argument}
For manuscript 1 (family 142; 2026-10-04), primary generators and norm equations reduce splitting to an odd-degree separator. Lifted ramification classes are filtered to keep extension degrees polynomial, and divisors at infinity distinguish roots; the analytic input bounds the auxiliary primes required by the finite algebraic algorithm.

\subsection{Scope relative to this problem}
Complete dense factorization over prime fields is covered, as are its quadratic applications. The bundle additionally asks about arbitrary finite fields, optimal randomized exponents, batched comparisons, and sparse-exponent root tests; a prime-field factorization theorem does not resolve every one of those tasks.

\subsection{Formal verification}
No Lean formalization is listed at this snapshot for manuscript 1 (family 142; 2026-10-04).

This site has not executed the Lean code. Described statement coverage is not an independent proof check.

\subsection{Status caveat}
The manuscripts were produced by an internal OpenAI model. OpenAI cautions that results without formalization can contain mistakes. This site has not independently checked the proofs; the label records the scope of OpenAI's claim.

\subsection{Sources}
\begin{itemize}
\item OpenAI, \emph{Deterministic Polynomial Factorization over Prime Fields}, 2026-10-04; citation key \texttt{OAI:Deterministic-Polynomial-Factorization-over-Prime-Fields-October-4-2026}. \href{https://github.com/openai/math/blob/adc7f1241b42e322a6451854ab7e4b4c146bf78a/preprints/Deterministic-Polynomial-Factorization-over-Prime-Fields-October-4-2026/Deterministic-Polynomial-Factorization-over-Prime-Fields.pdf}{PDF (GitHub)}; \href{https://github.com/openai/math/raw/adc7f1241b42e322a6451854ab7e4b4c146bf78a/preprints/Deterministic-Polynomial-Factorization-over-Prime-Fields-October-4-2026/Deterministic-Polynomial-Factorization-over-Prime-Fields.pdf}{PDF (direct)}; \href{https://github.com/openai/math/blob/adc7f1241b42e322a6451854ab7e4b4c146bf78a/preprints/Deterministic-Polynomial-Factorization-over-Prime-Fields-October-4-2026/README.md}{Manuscript page and citation}; \href{https://github.com/openai/math/blob/main/preprints/Deterministic-Polynomial-Factorization-over-Prime-Fields-October-4-2026/Deterministic-Polynomial-Factorization-over-Prime-Fields.pdf}{Latest version}; Statement source: \href{https://github.com/openai/math/blob/adc7f1241b42e322a6451854ab7e4b4c146bf78a/preprints/Deterministic-Polynomial-Factorization-over-Prime-Fields-October-4-2026/build/sections/00-introduction.tex\#L27}{thm:factorization}.
\item \href{https://github.com/openai/math/blob/adc7f1241b42e322a6451854ab7e4b4c146bf78a/CONTENTS.md}{OpenAI catalogue entry}.
\end{itemize}
\end{document}
